\begin{chapintro}
メタ解析の論文は、読者が同じ手順をたどり直せるように書きます。報告の指針 PRISMA 2020 の要点を確かめたうえで、降圧薬 A のレビューを例に方法と結果の段落を書いてみます。そのあと、他の人が書いたメタ解析を読むときに見るところをまとめ、第 1 部を締めくくります。本書で扱わない方法の紹介は発展の節に置いてあり、初読では飛ばせます。
\end{chapintro}

\section{PRISMA 2020 の要点}
\label{sec:08-1}

系統的レビューとメタ解析の報告には、PRISMA 2020 という指針を使います（Page 2021）。題名から資金源まで 27 の項目があり、それぞれに何を書くかが決められています。多くの雑誌は、投稿のときにこのチェックリストを埋めて、各項目を原稿の何ページに書いたかを示すよう求めます。抄録だけに使う 12 項目のチェックリストも別にあります。

27 項目の多くは、ここまでの章で行った作業をそのまま書けば満たせます。書き忘れやすいのは次の項目です。

\begin{description}
\item[題名] 系統的レビューであること（メタ解析を含むならそれも）が題名から分かるようにします。
\item[計画書と登録] PROSPERO の登録番号と、計画書から変えた点があればその内容と理由を書きます。変更を書かないと、後から解析を選んだのではないかと疑われます。
\item[検索] 使ったすべてのデータベース、登録簿、ウェブサイトについて、検索日と検索式の全文（使ったフィルターと制限を含む）を載せます。PRISMA 2020 の項目 7 がこれを求めています。全文は補足資料に置いて構いません。検索の報告をさらに詳しく書くための指針として、PRISMA-S（Rethlefsen 2021）があります。
\item[選別と抽出の手順] 何人が、独立に行ったか、食い違いをどう解決したかを書きます。
\item[効果の指標と統合の方法] リスク比を使ったこと、ランダム効果モデル、τ² の推定法、信頼区間の作り方を書きます。どれを選んだかで数字が変わるので、ソフトウェアとそのバージョンも書きます。
\item[報告バイアスと確実性の評価] 出版バイアスをどう調べたか、GRADE で確実性をどう判定したかを、方法の節に書きます。結果の節にだけ書くと、事前に決めた手順だったかどうかが読者に分かりません。
\item[個々の試験の結果] 統合値だけでなく、試験ごとの効果と信頼区間を示します。フォレストプロットがこれにあたります。
\item[データとコードの公開] 抽出したデータ、解析のコード、その他の資料を公開しているか、どこで手に入るかを書きます。
\end{description}

選別の経過は \kkref{sec:03-5} で描いた PRISMA のフロー図で示します。フロー図の数とチェックリストの数、本文の数がそろっているかは、投稿前に必ず突き合わせます。

\section{方法を書く}
\label{sec:08-2}

方法の節は、読者が同じ手順をたどれば同じ結果になるように書きます。降圧薬 A のレビューの方法を、架空の例として 2 段落で示します（以下の枠内はすべて架空のレビューです）。検索日と登録番号は架空のものです。

\begin{KKfbox9ao}[方法の例]
本レビューは PRISMA 2020 に従って報告し、計画書は事前に PROSPERO に登録した（CRD000000000）。高血圧の成人を対象に降圧薬 A とプラセボを比べ、脳卒中の発症を報告した RCT を対象とした。MEDLINE、Embase、CENTRAL、医中誌 Web、ClinicalTrials.gov と jRCT を、開始時から 2025 年 6 月 30 日まで、言語の制限なく検索した。題名と抄録の選別、本文の確認、データの抽出は 2 人が独立に行い、食い違いは話し合いで解決した。各試験のバイアスのリスクは RoB 2 で評価した。

効果の指標はリスク比とし、逆分散法によるランダム効果モデルで統合した。τ² は REML 法で推定し、95\% 予測区間を示した。異質性は I² と τ² で評価した。事前に決めたサブグループ解析として対象集団（高齢者、一般の成人）とバイアスのリスクの判定で分けた解析を、感度分析としてバイアスのリスクが高い試験を除いた解析、1 試験ずつ除いた解析、固定効果モデルと Hartung-Knapp 法による解析を行った。出版バイアスはファネルプロットと Egger 検定で調べ、固定効果モデルとランダム効果モデルの統合値を比べた。エビデンスの確実性は GRADE で判定した。解析には R 4.5.3 と meta パッケージ 8.5-0 を用いた。
\end{KKfbox9ao}

枠内は論文の文体（である調）で書いてあります。統計の段落には、効果の指標、モデル、τ² の推定法、信頼区間の作り方、異質性の指標、事前に決めた追加解析、ソフトウェアのバージョンを入れます。どれも数字を左右する選択なので、一つでも欠けると読者は結果を再現できません。例の中の逆分散法は重み付き平均の計算の仕方（\kkref{sec:05-3}）、REML は τ² の推定法、Hartung-Knapp 法は信頼区間の作り方で、後の二つは \kkref{sec:05-7} で説明しました。

\section{結果を書く}
\label{sec:08-3}

結果の節は、選別の経過、組み入れた試験の特徴、バイアスのリスク、統合の結果、追加解析、確実性の順に書くのが一般的です。統合の結果の段落を例に示します。

\begin{KKfbox9ao}[結果の例]
10 試験、21,420 人（脳卒中 1,304 件）を統合した。降圧薬 A はプラセボと比べて脳卒中を減らした（リスク比 0.73、95\% 信頼区間 0.61〜0.86、p < 0.001；図 2）。試験間の異質性は中程度で（I² = 51\%、τ² = 0.034）、95\% 予測区間は 0.46〜1.15 であった。対象集団で分けると、リスク比は高齢者で 0.65（0.46〜0.91）、一般の成人で 0.82（0.71〜0.95）であったが、サブグループ間の差は有意ではなかった（p = 0.20）。バイアスのリスクが高い 1 試験を除いた解析（0.74、0.62〜0.88）、1 試験ずつ除いた解析（0.70〜0.77）、Hartung-Knapp 法による解析（0.73、0.59〜0.89）でも結果は変わらなかった。Egger 検定でファネルプロットの非対称は検出されなかった（p = 0.27）。確実性は、非一貫性のため中と判定した（表 3）。
\end{KKfbox9ao}

数字の書き方にはいくつかの約束があります。点推定には必ず信頼区間を添え、p 値だけを書くことはしません。リスク比は小数 2 桁、p 値は 2 桁程度にし、0.001 より小さければ p < 0.001 とします。サブグループ解析は、群ごとの値だけでなく群の間の差の検定を書きます。群ごとの値だけを並べると、片方だけ有意だったことを差があったと読み違えられやすいからです（\kkref{sec:06-4}）。

図表は、PRISMA のフロー図、試験の特徴の表、バイアスのリスクの図、フォレストプロット、結果の要約表をそろえるのが基本です。フォレストプロットの試験の順番、表の試験の順番、本文で試験を挙げる順番をそろえておくと、読者が照らし合わせやすくなります。

\section{投稿前の確認と公開}
\label{sec:08-4}

投稿の前に、次の点を確かめます。

\begin{description}
\item[数の一致] 抄録、本文、表、図、フロー図の数（試験数、人数、事象数、効果の推定値）が一致しているか。修正を繰り返すうちに、一か所だけ古い数字が残ることがよくあります。
\item[計画書との一致] 計画書に書いた解析をすべて報告しているか。変えた点は理由とともに書いてあるか。
\item[チェックリスト] PRISMA 2020 の各項目について、書いたページを埋められるか。
\item[再現] 解析のコードを最初から実行し直して、論文の数字が再現されるか。
\end{description}

抽出したデータと解析のコードは、補足資料に付けるか、OSF や GitHub のような公開の場所に置きます。メタ解析のデータは発表済みの論文から取った試験単位の集計値なので、個人を特定する情報は含まれません。公開しておけば、読者が数字を確かめられ、次にレビューを更新する人がそこから作業を始められます。本書の R スクリプトのように、データの読み込みから図の書き出しまでを一つのスクリプトにまとめておくと、再現が容易になります。

\section{メタ解析の論文を読むときに見るところ}
\label{sec:08-5}

自分でメタ解析をしなくても、検査技師として診療ガイドラインや総説の根拠になったメタ解析を読む機会はあります。結論を受け取る前に、次の点を確かめます。

\begin{description}
\item[問いと対象] PICO がはっきりしているか。対象、介入、比較、アウトカムが、自分の知りたい問いと合っているか。
\item[計画書の登録] 事前に登録されているか。登録された計画と報告された解析が一致しているか。
\item[検索の範囲] 複数のデータベースと試験登録を検索しているか。言語や発表形式で絞っていないか。
\item[選別と抽出] 2 人が独立に行ったか。除外した理由が示されているか。
\item[個々の試験の質] バイアスのリスクを評価しているか。その評価が結論に反映されているか。
\item[統合の妥当性] 統合した試験の対象や介入は似ているか。異質性が大きいのに一つの数字だけで結論していないか。予測区間は示されているか。
\item[サブグループ解析] 事前に決めたものか。群の間の差を検定しているか。数多くのサブグループのうち一つだけを強調していないか。
\item[出版バイアス] 調べているか。試験数が少ないのに Egger 検定の結果だけで「ない」としていないか。
\item[確実性] GRADE などで判定しているか。相対効果だけでなく絶対効果が示されているか。
\end{description}

全部の項目を満たす論文は多くありません。どこが弱いかを把握したうえで結論を読めば、その結論をどこまで診療や検査の判断に使えるかを見積もれます。

\section{\hatten この先に進むときに}
\label{sec:08-6}

この冊子で扱ったのは、2 つの介入を直接比べた試験を、公表された集計値で統合する方法です。次の 4 つは本書では扱いません。必要になったときの手がかりとして、Cochrane ハンドブックの該当箇所を挙げます。

\begin{description}
\item[ネットワークメタ解析（第 11 章）] 3 つ以上の介入を、直接比べた試験と間接的な比較を組み合わせて一度に比べる方法です。間接比較が成り立つには、試験どうしが効果に影響する要因について似ている必要があり（推移性、11.1.3）、これを確かめる作業が解析の前提になります。
\item[個人データのメタ解析（第 26 章）] 試験ごとの集計値ではなく、参加者一人ひとりのデータを集めて統合する方法です（26.1.1）。未報告のアウトカムを使え、解析方法を試験間でそろえられ、年齢や重症度によって効果が違うかを調べる力も高くなりますが、データの提供を受ける手間がかかります。
\item[リビングシステマティックレビュー（22.2.3）] 新しい証拠が出るたびに、できるだけ早く取り込んで更新し続ける系統的レビューです。新しい試験を見つけ続ける監視の仕組み（22.2）と組み合わせて運用します。
\item[ベイズ流のメタ解析（10.13）] 事前の情報を事前分布として置き、データと合わせて事後分布を求めます。結果は点推定と信用区間で示され、事前分布の選び方に結果が左右されるので、感度分析が勧められています。試験数が少なく τ² の推定が不安定なときの選択肢としても挙げられています。
\end{description}

この 4 つはどれも、第 1 部と同じく介入の効果を 1 つの数で統合する流れの中で使う方法です。

\section{第 1 部のまとめ}
\label{sec:08-7}

第 1 部では、降圧薬 A の 10 試験を使って、介入のメタ解析を問いの設定から報告まで一通りたどりました。問いを PICO で定めて計画書を登録し（\kkref{sec:02-1}、\kkref{sec:02-4}）、試験を検索して選別し（\kkref{sec:03-2}、\kkref{sec:03-3}）、データを取り出して偏りを評価しました（\kkref{sec:04-2}、\kkref{sec:04-3}）。R の meta パッケージでリスク比を統合し（\kkref{sec:05-6}）、異質性を調べ（\kkref{sec:06-1}）、出版バイアス（\kkref{sec:07-1}）、感度分析（\kkref{sec:07-6}）、GRADE（\kkref{sec:07-8}）で結果の確かさを評価しました。

結果は、リスク比 0.73（0.61〜0.86）、確実性は中です。降圧薬 A は脳卒中を減らすと考えられますが、予測区間（\kkref{sec:06-3}）が 1 をまたいでおり、効果の大きさは試験によって違います。この結論は数字一つではなく、試験の集め方（\kkref{sec:03-1}）、統合の方法（\kkref{sec:05-4}）、異質性（\kkref{sec:06-1}）と偏り（\kkref{sec:07-1}）の評価を合わせて初めて言えるものです。

第 2 部では、プロカルシトニンの 10 研究を使って、検査の診断精度のメタ解析を行います。検索、選別、抽出、報告の流れは第 1 部と同じですが、検査の診断精度は感度と特異度（\kkref{sec:09-2}）という 2 つの数が互いに関係して動くので、1 つの数で統合する流れそのものが変わります。
